%% 19_dstheory.tex -- Chapter 19: $\Delta\Sigma$ Modulation: Theory
%% Status: OUTLINE (build stage 2). Replace the outline with the full chapter when it is written.
\chapter{$\Delta\Sigma$ Modulation: Theory}\label{ch:dstheory}
\chapterlead{Oversample, quantise coarsely and push the quantisation noise out of the band: noise shaping turns a 1-bit quantiser into a 16-bit converter.}

\begin{specnote}[Chapter status: outline]
This chapter is an outline. It will be written in build stage 2 (the Python-only converter core: Chapters 13, 14 and 19 with Labs 6 and 9). The topics, the design question and the labs below are fixed; the derivations, worked examples and figures are still to come.
\end{specnote}

\begin{objectives}
\begin{itemize}
  \item derive the noise transfer function (NTF) and signal transfer function (STF) of first- and second-order loops;
  \item predict SQNR from loop order and oversampling ratio (OSR);
  \item check stability with the Lee criterion and understand idle tones and dither;
  \item synthesise and realise an NTF with \code{synthesizeNTF} and \code{realizeNTF};
  \item know when linearised analysis fails and simulation is needed.
\end{itemize}
\end{objectives}

\begin{insight}[Design question]
Which loop order and OSR give 16 bits in the ECG/EEG band at the lowest clock rate?
\end{insight}

\section{Noise shaping}\label{sec:ch19-1}
The linear model of a $\Delta\Sigma$ loop; $\mathrm{NTF}(z) = (1-z^{-1})^L$ and its high-pass shape.

\begin{keyeq}[Peak SQNR of an order-$L$ modulator]
\begin{equation}
  \mathrm{SQNR}_{L} \approx 6.02N + 1.76 - 10\log_{10}\!\frac{\pi^{2L}}{2L+1} + (20L+10)\log_{10}\mathrm{OSR}\ \mathrm{dB}
  \label{eq:dssqnr}
\end{equation}
\end{keyeq}

\section{SQNR versus order and OSR}\label{sec:ch19-2}
Each doubling of OSR buys $6L + 3$\,dB; diminishing returns and the trade-off with stability.

\section{Stability, idle tones and dither}\label{sec:ch19-3}
The Lee criterion $\lVert \mathrm{NTF}\rVert_\infty < 1.5$ for 1-bit loops; tones with DC inputs, which matter for ECG baselines.

\section{NTF synthesis}\label{sec:ch19-4}
Using python-deltasigma~\cite{pydeltasigma}, the Python port of Schreier's toolbox~\cite{schreier}, to place zeros and poles and map them to a loop topology.

\begin{tryit}[Labs for this chapter]
Lab~9 (\cref{sec:lab9}) synthesises first- to third-order NTFs, simulates SQNR versus OSR and the incremental mode, and designs a fixed-point CIC decimator.
\end{tryit}

\paragraph{Read with this chapter} Schreier, Temes and Pavan~\cite{schreier}, python-deltasigma~\cite{pydeltasigma}.
